\documentclass[../report.tex]{subfiles}
\begin{document}

\section{Chia datapath thành năm tầng}

Pipeline cắt đường tổ hợp của CPU đơn chu kỳ thành năm đoạn IF, ID, EX, MEM, WB và đặt một thanh ghi pipeline giữa
hai đoạn liên tiếp. Mỗi chu kỳ, năm lệnh khác nhau nằm ở năm tầng. Chu kỳ clock chỉ cần đủ cho tầng chậm nhất cộng
với thời gian của thanh ghi pipeline. Với giả định trễ ở mục~\ref{sec:critical} và 20~ps cho thanh ghi pipeline,
chu kỳ pipeline là $\max(200, 100, 200, 200, 100) + 20 = \resTPipe$~ps, so với \resTSingle~ps của bản đơn chu kỳ.
Nếu mọi lệnh đều có CPI bằng 1, tăng tốc lý tưởng là $\resTSingle/\resTPipe \approx \resIdealSpeedup$ lần.

\begin{figure}[H]
\centering
\includegraphics[width=\textwidth]{RV32I_Pipeline.png}
\caption{Mạch top-level của CPU pipeline. Pipeline Monitor ở góc dưới hiển thị PC và lệnh ở từng tầng cùng các
tín hiệu \sig{Stall}, \sig{LoadUse}, \sig{E\_Redirect}, \sig{FwdA}, \sig{FwdB}}\label{fig:pl}
\end{figure}

\section{Thanh ghi pipeline}

Mỗi thanh ghi pipeline giữ mọi thứ các tầng sau cần về một lệnh: dữ liệu, số hiệu thanh ghi và các tín hiệu điều
khiển. Tín hiệu điều khiển được sinh một lần ở ID rồi đi theo lệnh qua từng tầng. Mỗi trường gồm một MUX
\sig{FLUSH} ? 0 : D và một Register có enable = \sig{EN} OR \sig{FLUSH}. Khi \sig{EN} = 0 thanh ghi giữ giá trị
(stall). Khi \sig{FLUSH} = 1 thanh ghi nạp toàn 0; do lệnh 0 có mọi tín hiệu điều khiển bằng 0, lệnh đó thành
bubble và không ghi gì. Flush được làm đồng bộ qua MUX vì chân Clear của Register trong Logisim là bất đồng bộ.

\begin{table}[H]
\centering\small
\caption{Các trường của thanh ghi pipeline}\label{tab:pregs}
\begin{tabularx}{\textwidth}{@{}lX@{}}
\toprule
Thanh ghi & Trường \\
\midrule
IF/ID & valid, pc, pc4, instr \\
ID/EX & valid, pc, pc4, instr, rs1d, rs2d, imm, rs1, rs2, rd, funct3, ALUCtrl, ASel, BSel, RegWrite, MemRead, MemWrite, WBSel, Branch, Jal, Jalr, Halt \\
EX/MEM & valid, pc, instr, exres, sdata, rd, funct3, RegWrite, MemRead, MemWrite, MemToReg, Halt \\
MEM/WB & valid, pc, instr, exres, mdata, rd, RegWrite, MemToReg, Halt \\
\bottomrule
\end{tabularx}
\end{table}

Trường \sig{exres} là PC + 4 với JAL/JALR và kết quả ALU với các lệnh khác. Nó vừa là địa chỉ bộ nhớ, vừa là
giá trị để forward, vừa là giá trị ghi lại, nên từ tầng MEM trở đi chỉ cần MUX 2:1 theo \sig{MemToReg}. Trường
\sig{valid} bằng 0 với bubble và dùng cho bộ đếm \sig{instret}. Các trường \sig{pc}, \sig{instr} ở mọi tầng chỉ
để hiển thị trên Pipeline Monitor.

\section{Ba loại hazard}

Structural hazard không xảy ra vì bộ nhớ lệnh và bộ nhớ dữ liệu tách riêng, register file có hai cổng đọc và một
cổng ghi. Data hazard xảy ra khi một lệnh cần kết quả của lệnh trước mà lệnh đó chưa ghi vào register file.
Control hazard xảy ra vì pipeline đã nạp tiếp hai lệnh sau một lệnh nhảy trước khi biết lệnh đó có nhảy hay không.

\subsection{\texorpdfstring{Forwarding và bypass WB$\to$ID}{Forwarding và bypass WB-ID}}\label{sec:fwd}

Kết quả của một lệnh ALU đã có ở cuối tầng EX, chỉ chưa được ghi vào register file. \sig{ForwardUnit} so sánh
rs1, rs2 của lệnh ở EX với rd của lệnh ở MEM và WB. \sig{FwdA} = 01 (lấy từ EX/MEM) khi EX/MEM.RegWrite = 1,
EX/MEM.rd $\neq$ 0 và EX/MEM.rd = ID/EX.rs1. \sig{FwdA} = 10 (lấy từ MEM/WB) khi điều kiện tương tự với MEM/WB đúng
còn điều kiện 01 sai. Các trường hợp khác cho \sig{FwdA} = 00, dùng giá trị đọc ở ID. \sig{FwdB} làm tương tự với rs2.

Điều kiện rd $\neq$ 0 ngăn lệnh \texttt{add x0, \ldots} forward một giá trị khác 0 cho lệnh sau đọc x0. EX/MEM được
ưu tiên vì khi hai lệnh liên tiếp cùng ghi một thanh ghi, giá trị mới nhất thuộc lệnh gần hơn. MUX forwarding đặt
trước MUX \sig{ASel}/\sig{BSel}, nên giá trị đã forward dùng được cho ALU, cho \sig{BranchUnit} và làm dữ liệu
store.

Còn trường hợp lệnh cách 3 vị trí: lệnh ở WB ghi register file ở cạnh lên cuối chu kỳ, trong khi lệnh ở ID đọc
ngay trong chu kỳ đó và nhận giá trị cũ. \sig{IDBypass} đặt sau cổng đọc của register file: nếu W\_RegWrite = 1,
W\_rd $\neq$ 0 và W\_rd trùng rs thì đưa thẳng W\_data vào ID. Cách này tương đương với quy ước ghi nửa đầu, đọc
nửa sau chu kỳ trong sách \cite[ch.~4]{patterson}, và giữ \sig{RegisterFile} dùng được cho CPU đơn chu kỳ.

\subsection{Load-use: stall một chu kỳ}

Dữ liệu của lệnh load chỉ có ở cuối tầng MEM, còn lệnh ngay sau cần nó ở đầu tầng EX trong cùng chu kỳ, nên
forwarding không đủ. \sig{HazardUnit} phát hiện trường hợp này:

\begin{align*}
\mathit{LoadUse} = {}& \mathit{E\_MemRead} \land (E\_rd \neq 0) \\
&\land \big((E\_rd = rs1 \land \mathit{UseRs1}) \lor (E\_rd = rs2 \land \mathit{UseRs2})\big)
\end{align*}

Khi \sig{LoadUse} = 1, \sig{PCWrite} = 0 và \sig{IFID\_Write} = 0 giữ nguyên PC và IF/ID, còn \sig{IDEX\_Flush} = 1
chèn một bubble vào EX. Chu kỳ sau, lệnh phụ thuộc lên EX và nhận dữ liệu load qua forwarding từ MEM/WB.

\subsection{Control hazard: dự đoán không nhảy và flush}

PL luôn nạp PC + 4. Lệnh rẽ nhánh và nhảy được quyết định ở EX: \sig{E\_Redirect} = \sig{Taken} $\lor$ \sig{Jal}
$\lor$ \sig{Jalr}. Khi \sig{E\_Redirect} = 1, PC nhận địa chỉ đích và hai lệnh nạp sai ở IF/ID và ID/EX bị flush,
mất 2 chu kỳ. Nếu cùng chu kỳ có cả stall và redirect thì redirect thắng, vì lệnh đang bị stall nằm trên nhánh sai.
Các phương trình của \sig{HazardUnit} là:

\[
\mathit{Stall} = \mathit{LoadUse} \lor \mathit{D\_Halt},\quad
\mathit{PCWrite} = \mathit{IFID\_Write} = \lnot\mathit{Stall} \lor \mathit{Redirect},
\]
\[
\mathit{IFID\_Flush} = \mathit{Redirect},\quad \mathit{IDEX\_Flush} = \mathit{Redirect} \lor \mathit{LoadUse}.
\]

EBREAK ở ID làm \sig{Stall} = 1 để ngừng nạp lệnh mới. Chân \sig{halt} lấy từ trường \sig{Halt} của MEM/WB, nên
CPU chỉ dừng khi EBREAK tới WB, lúc mọi lệnh trước đó đã ghi xong.

Từ các cơ chế trên, một chương trình $N$ lệnh chạy trên PL mất
$N + 4 + (\text{số lần load-use}) + 2 \times (\text{số lần nhảy})$ chu kỳ, trong đó 4 là số chu kỳ lấp đầy pipeline.
Bộ test so đúng con số này với bộ đếm \sig{cycles} của Logisim trên mọi chương trình (chương~\ref{ch:test}).

\section{Bản NoHazard và cách chèn NOP}\label{sec:nh}

Để đo lợi ích của phần cứng xử lý hazard, nhóm giữ một bản pipeline có hai thay đổi so với PL: chân chọn của MUX
forwarding nối hằng 0, và đầu vào \sig{E\_MemRead} của \sig{HazardUnit} nối hằng 0. Flush khi nhảy và
\sig{IDBypass} vẫn giữ. Chạy chương trình bình thường trên bản này thì sai: cả 12 chương trình test đều FAIL
(mục~\ref{sec:testres}).

Để bản NH chạy đúng, assembler có chế độ \texttt{--nohazard} chèn \texttt{nop} theo hai quy tắc. Lệnh đọc rd của
lệnh ngay trước được chèn 2 \texttt{nop}; lệnh đọc rd của lệnh trước đó 2 vị trí được chèn 1 \texttt{nop}. Sau đó
khoảng cách giữa hai lệnh là 3 và \sig{IDBypass} lo phần còn lại. Bộ đếm khoảng cách được đặt lại sau
\texttt{jal} và \texttt{jalr}, vì lệnh đích chỉ đến được qua một lần nhảy và mỗi lần nhảy đã flush 2 chu kỳ.

Bảng~\ref{tab:diag} so sánh cùng bảy lệnh của \file{t06\_hazard.s} trên hai bản. Ô màu đỏ là chu kỳ lệnh bị stall
ở ID; bảng sinh tự động từ golden model.

\begin{table}[H]
\centering\scriptsize
\caption{Giản đồ pipeline của 7 lệnh trong \file{t06\_hazard.s}}\label{tab:diag}
\setlength{\tabcolsep}{1.6pt}
\begin{subtable}{\textwidth}\centering
\caption{PL: forwarding và stall bằng phần cứng}
\begin{adjustbox}{max width=\textwidth}
\input{generated/t06_pl.tex}
\end{adjustbox}
\end{subtable}

\vspace{0.8em}
\begin{subtable}{\textwidth}\centering
\caption{NH: không forwarding, assembler chèn \texttt{nop}}
\begin{adjustbox}{max width=\textwidth}
\input{generated/t06_nh.tex}
\end{adjustbox}
\end{subtable}
\end{table}

Trên PL, ba lệnh dùng ngay kết quả của \texttt{lw} mỗi lệnh chờ 1 chu kỳ, các phụ thuộc khác được forward. Lệnh
\texttt{beq} nhảy ở EX nên hai lệnh phía sau bị flush. Bảy lệnh chiếm 14 chu kỳ từ lúc \texttt{sw} vào IF đến lúc
\texttt{beq} rời WB. Trên NH, mỗi phụ thuộc gần phải chờ bằng \texttt{nop}, kể cả phụ thuộc mà PL forward được không
mất chu kỳ nào, nên cùng đoạn này chiếm 17 chu kỳ.

\end{document}
